\documentclass[12pt,a4paper,oneside]{report}

\usepackage{geometry}
\geometry{
    top=1.5in,
    bottom=1.0in,
    left=1.5in,
    right=1.0in
}

\usepackage{fontspec}
\usepackage{xunicode}
\usepackage{xltxtra}

\defaultfontfeatures{Mapping=tex-text}
\setmainfont{THSarabunNew.ttf}[
    Path = /Users/chewathassana/Library/Fonts/,
    BoldFont = {THSarabun Bold.ttf},
    ItalicFont = {THSarabun Italic.ttf},
    BoldItalicFont = {THSarabun BoldItalic.ttf}
]

\XeTeXlinebreaklocale "th_TH"
\XeTeXlinebreakskip = 0pt plus 1pt

\usepackage{amsmath,amssymb,amsfonts}
\usepackage{graphicx}
\usepackage{xcolor}
\usepackage{booktabs}
\usepackage{tabularx}
\newcolumntype{Y}{>{\raggedright\arraybackslash}X}

\usepackage{tcolorbox}
\tcbuselibrary{skins,breakable}
\usepackage{tikz}
\usetikzlibrary{shapes,arrows,positioning,calc,fit,backgrounds}
\usepackage{titlesec}
\usepackage{caption}

% กำหนดค่าสีตามอัตลักษณ์มหาวิทยาลัยราชภัฏรำไพพรรณีและธีมวิจัย
\definecolor{rbruNavy}{RGB}{15, 30, 55}
\definecolor{rbruGold}{RGB}{205, 155, 30}
\definecolor{rbruSlate}{RGB}{70, 85, 105}
\definecolor{rbruGreen}{RGB}{20, 130, 50}
\definecolor{rbruCodeBg}{RGB}{245, 248, 252}
\definecolor{isoilCyan}{RGB}{0, 180, 220}

% การตั้งค่าคำบรรยายภาพและตารางตามมาตรฐาน RBRU (ห้ามมีเครื่องหมายทวิภาค)
\DeclareCaptionLabelSeparator{thaisep}{\quad}
\captionsetup{labelsep=thaisep}
\renewcommand{\figurename}{\mbox{ภาพที่}}
\renewcommand{\tablename}{\mbox{ตารางที่}}

\AtBeginDocument{
    \gdef\figurename{\mbox{ภาพที่}}
    \gdef\tablename{\mbox{ตารางที่}}
    \captionsetup[figure]{name={\mbox{ภาพที่}}, font={small}}
    \captionsetup[table]{name={\mbox{ตารางที่}}, font={small}, position=top}
}

% สไตล์หัวบทวิชาการหรูหรา
\newcommand{\rbruchapterbanner}[1]{%
    \begin{tcolorbox}[
        enhanced,
        colback=rbruNavy,
        colframe=rbruGold,
        arc=4mm,
        boxrule=1.5pt,
        left=6mm, right=6mm, top=6mm, bottom=6mm
    ]
    \raggedright
    {\color{rbruGold}\fontsize{14}{16}\selectfont\textbf{คู่มือปฏิบัติการวิจัยและเทคโนโลยีปัญญาประดิษฐ์ฝังตัว}}\par\vspace{2mm}
    {\color{white}\fontsize{22}{26}\selectfont\textbf{บทที่ \thechapter\quad #1}}
    \end{tcolorbox}
}

\titleformat{\chapter}[block]{\normalfont}{}{0pt}{\rbruchapterbanner}
\titlespacing*{\chapter}{0pt}{-10pt}{20pt}

\titleformat{\section}[block]{\normalfont\fontsize{16}{20}\selectfont\bfseries\color{rbruNavy}}{\thesection\quad}{0pt}{}
\titleformat{\subsection}[block]{\normalfont\fontsize{14}{18}\selectfont\bfseries\color{rbruSlate}}{\thesubsection\quad}{0pt}{}

\setlength{\parindent}{1.0cm}
\setlength{\parskip}{4pt}

\begin{document}

% ---------------------- หน้าปกเอกสารทางวิชาการ ----------------------
\begin{titlepage}
\begin{tcolorbox}[
    enhanced,
    colback=white,
    colframe=rbruNavy,
    boxrule=2pt,
    arc=0mm,
    height=\textheight,
    interior style={left color=white, right color=rbruCodeBg}
]
\centering
\vspace*{0.8cm}
{\color{rbruGold}\rule{0.85\textwidth}{3pt}}\par\vspace{5mm}
{\fontsize{26}{30}\selectfont\textbf{\color{rbruNavy}คู่มือฉบับสมบูรณ์สำหรับผู้เริ่มต้น}}\par\vspace{3mm}
{\fontsize{20}{24}\selectfont\textbf{\color{isoilCyan}ระบบ iSoil pH xAI}}\par\vspace{2mm}
{\fontsize{16}{20}\selectfont\textbf{\color{rbruSlate}เครื่องวัดความเป็นกรด-ด่างของดินภาคสนามความแม่นยำสูง}}\par
{\fontsize{15}{19}\selectfont\textbf{\color{rbruGreen}ด้วยการชดเชยความคลาดเคลื่อนโดยปัญญาประดิษฐ์ฝังตัว (TinyML)}}\par\vspace{3mm}
{\fontsize{12}{15}\selectfont\textit{Development of a High-Accuracy Field-Portable Soil pH Meter Prototype\\with Explainable TinyML AI Error Compensation}}\par\vspace{5mm}
{\color{rbruGold}\rule{0.85\textwidth}{1.5pt}}\par\vspace{1.0cm}

\begin{tikzpicture}[scale=0.85]
    \draw[fill=rbruNavy!10, draw=rbruNavy, line width=1.5pt, rounded corners=6mm] (0,0) rectangle (7.5,4.8);
    \node at (3.75,4.1) {\textbf{\color{rbruNavy}iSoil pH xAI [ANN 2-4-1 MLP]}};
    \draw[fill=white, draw=rbruGold, line width=1pt, rounded corners=3mm] (0.8,0.8) rectangle (6.7,3.5);
    \node at (3.75,2.7) {\fontsize{28}{32}\selectfont\textbf{\color{rbruGreen}pH 6.85}};
    \node at (3.75,1.8) {\fontsize{12}{14}\selectfont\textbf{\color{rbruGold}RAW CELL: 1.661 V}};
    \node at (3.75,1.2) {\fontsize{10}{12}\selectfont\textbf{\color{rbruSlate}OPTIMAL DURIAN / pH 5.5-6.5}};
    \draw[line width=1.5pt, color=rbruNavy] (7.5,2.4) -- (9.0,2.4);
    \draw[fill=rbruSlate!20, draw=rbruSlate, line width=1pt] (9.0,1.7) rectangle (11.0,3.2);
    \node at (10.0,2.4) {\fontsize{10}{12}\selectfont\textbf{BNC PROBE}};
\end{tikzpicture}

\vfill

{\fontsize{14}{17}\selectfont\textbf{คณะผู้วิจัย}}\par\vspace{2mm}
{\fontsize{13}{16}\selectfont
อาจารย์ธนพัฒน์ ถิระวุฒิ\\
ผู้ช่วยศาสตราจารย์ ดร.ชีวะ ทัศนา\\
ผู้ช่วยศาสตราจารย์ ดร.นันทพร มูลรังษี
}\par\vspace{6mm}

{\fontsize{12}{15}\selectfont
โครงการวิจัยภายใต้กองทุนวิจัย มหาวิทยาลัยราชภัฏรำไพพรรณี\\
ประจำปีงบประมาณ พ.ศ. 2569\\
คณะวิทยาศาสตร์และเทคโนโลยี มหาวิทยาลัยราชภัฏรำไพพรรณี
}
\vspace*{0.6cm}
\end{tcolorbox}
\end{titlepage}

% ---------------------- บทที่ 1 ----------------------
\chapter{บทนำและหลักการเกษตรฟิสิกส์}

ความเป็นกรด-ด่างของดิน (Soil pH) เป็นดัชนีชี้วัดทางเคมีฟิสิกส์ที่กำหนดความสามารถในการละลายและการดูดซึมธาตุอาหารของพืช โดยเฉพาะอย่างยิ่งในพื้นที่เกษตรกรรมของจังหวัดจันทบุรีและจังหวัดตราด ซึ่งเป็นแหล่งผลิตพืชเศรษฐกิจสำคัญของประเทศ ได้แก่ ทุเรียน (\textit{Durio zibethinus} Murr.) ดินส่วนใหญ่ในเขตนี้มักเผชิญภาวะดินกรดรุนแรงเนื่องจากมีปริมาณน้ำฝนสะสมสูงตลอดทั้งปี

\begin{figure}[htbp]
\centering
\begin{tikzpicture}[scale=0.85, auto, >=stealth]
    \node[draw=rbruNavy, fill=rbruCodeBg, rounded corners, inner sep=3mm, align=center] (input) at (0,0) {\textbf{สารละลายดิน}\\\fontsize{9}{11}\selectfont (Soil Slurry 1:1)};
    \node[draw=rbruGold, fill=white, rounded corners, inner sep=3mm, align=center] (sensor) at (3.8,0) {\textbf{หัววัดแก้ว BNC}\\\fontsize{9}{11}\selectfont (Glass Electrode)};
    \node[draw=rbruSlate, fill=white, rounded corners, inner sep=3mm, align=center] (amp) at (7.8,0) {\textbf{pH-4502C}\\\fontsize{9}{11}\selectfont (Offset 1.650V)};
    \node[draw=rbruGreen, fill=rbruGreen!10, rounded corners, inner sep=3mm, align=center] (mcu) at (11.8,0) {\textbf{Wio Terminal}\\\fontsize{9}{11}\selectfont (TinyML ANN)};

    \draw[->, line width=1.3pt] (input) -- node[above]{\fontsize{8}{10}\selectfont ไอออน $H^+$} (sensor);
    \draw[->, line width=1.3pt] (sensor) -- node[above]{\fontsize{8}{10}\selectfont mV ศักย์} (amp);
    \draw[->, line width=1.3pt] (amp) -- node[above]{\fontsize{8}{10}\selectfont 0 - 3.3V} (mcu);
\end{tikzpicture}
\caption{ผังการทำงานเชิงฟิสิกส์ของการแปลงความเข้มข้นไฮโดรเจนไอออนสู่การคำนวณด้วยปัญญาประดิษฐ์}
\label{fig:system_overview}
\end{figure}

เมื่อดินมีค่า pH ต่ำกว่า 5.0 ไอออนของโลหะหนักโดยเฉพาะอะลูมิเนียม ($Al^{3+}$) จะละลายออกมาเป็นพิษต่อระบบราก ขัดขวางการเจริญเติบโต และทำให้ผลผลิตตกต่ำ ในทางกลับกัน ค่า pH ที่เหมาะสมที่สุดสำหรับต้นทุเรียนจะอยู่ในช่วง 5.5 ถึง 6.5 ซึ่งเป็นระดับที่ธาตุอาหารหลักและธาตุอาหารรองละลายออกมาได้อย่างสมดุล

% ---------------------- บทที่ 2 ----------------------
\chapter{รากฐานทางทฤษฎีและสมการเนิร์นสต์}

การตรวจวัดค่าความเป็นกรด-ด่างด้วยวิธีโพเทนชิออเมตริก (Potentiometric Method) อาศัยความต่างศักย์ไฟฟ้าระหว่างขั้วไฟฟ้าบ่งชี้เยื่อแก้ว (Glass Indicator Electrode) และขั้วไฟฟ้าอ้างอิง (Reference Electrode) ซึ่งอธิบายได้ด้วยสมการเนิร์นสต์ (Nernst Equation) ดังนี้

\begin{equation}
E = E^0 - \frac{2.303 R T}{n F} \cdot \text{pH}
\end{equation}

เมื่อกำหนดให้
\begin{itemize}
    \item $E$ คือ ศักย์ไฟฟ้าที่อ่านได้จากเซลล์เคมีไฟฟ้า มีหน่วยเป็นโวลต์ (V)
    \item $E^0$ คือ ศักย์ไฟฟ้ามาตรฐานของเซลล์ มีหน่วยเป็นโวลต์ (V)
    \item $R$ คือ ค่าคงที่สากลของแก๊ส ($8.314\text{ J}\cdot\text{mol}^{-1}\cdot\text{K}^{-1}$)
    \item $T$ คือ อุณหภูมิสัมบูรณ์ของสารละลาย มีหน่วยเป็นเคลวิน (K)
    \item $n$ คือ จำนวนอิเล็กตรอนในปฏิกิริยา ซึ่งมีค่าเท่ากับ 1 สำหรับไฮโดรเจนไอออน
    \item $F$ คือ ค่าคงที่ฟาราเดย์ ($96,485\text{ C}\cdot\text{mol}^{-1}$)
\end{itemize}

\begin{table}[htbp]
\centering
\caption{ค่าความชันเนิร์นสเตียนตามทฤษฎีที่ระดับอุณหภูมิต่างๆ}
\label{tab:nernst_slope}
\begin{tabularx}{\textwidth}{cccY}
\toprule
\textbf{อุณหภูมิ ($^\circ\text{C}$)} & \textbf{อุณหภูมิสัมบูรณ์ (K)} & \textbf{ความชัน ($S$, mV/pH)} & \textbf{ลักษณะการนำไปใช้งาน} \\
\midrule
20.0 & 293.15 & 58.17 & สภาวะห้องทดลองควบคุมอุณหภูมิ \\
25.0 & 298.15 & 59.16 & อุณหภูมิมาตรฐานสากลอ้างอิง \\
30.0 & 303.15 & 60.15 & สภาวะแปลงเกษตรภาคสนามช่วงเช้า \\
35.0 & 308.15 & 61.14 & สภาวะแปลงเกษตรภาคสนามกลางแจ้ง \\
40.0 & 313.15 & 62.13 & สภาวะแดดจัดในดินภาคตะวันออก \\
45.0 & 318.15 & 63.13 & สภาวะทดสอบความทนทานขีดสุด \\
\bottomrule
\end{tabularx}
\end{table}

จากตารางจะเห็นได้ชัดเจนว่า หากไม่มีการชดเชยผลของอุณหภูมิ ($T$) และความไม่เป็นเชิงเส้นของเยื่อแก้ว การวัดค่าภาคสนามจะคลาดเคลื่อนสูงมาก โครงการวิจัยนี้จึงนำแบบจำลองโครงข่ายประสาทเทียม (ANN) มาทำหน้าที่เป็นฟังก์ชันผกผันเพื่อชดเชยความคลาดเคลื่อนดังกล่าว

% ---------------------- บทที่ 3 ----------------------
\chapter{ฮาร์ดแวร์ ผังวงจร และการสร้างต้นแบบ}

ต้นแบบเครื่องวัดค่าความเป็นกรด-ด่างของดินภาคสนามประกอบด้วย 3 ส่วนหลัก ได้แก่ บอร์ดไมโครคอนโทรลเลอร์ Seeed Studio Wio Terminal บอร์ดทดลอง INEX NX-WIO และชุดโมดูลวัดค่า pH-4502C พร้อมโพรบแก้ว

\begin{table}[htbp]
\centering
\caption{การต่อสายวงจรระหว่างเซนเซอร์ pH-4502C กับบอร์ด INEX NX-WIO}
\label{tab:wiring}
\begin{tabularx}{\textwidth}{p{3.5cm}p{4.0cm}Y}
\toprule
\textbf{ขาของโมดูล pH-4502C} & \textbf{ตำแหน่งบนบอร์ด NX-WIO} & \textbf{หน้าที่และความสำคัญเชิงวิศวกรรม} \\
\midrule
V+ (VCC) & ขั้วไฟ +Vs (max 5V) & จ่ายไฟเลี้ยง 5.0V ให้วงจรขยายสัญญาณ Op-Amp \\
G (GND)  & ขั้วกราวด์ GND       & กราวด์ร่วมของระบบเพื่อความเสถียรของสัญญาณ \\
Po (Analog Out) & ช่องเสียบเลข [27] (ขา A0) & สัญญาณแรงดันแอนะล็อกเชื่อมต่อเข้าสู่ ADC 12 บิต \\
ขั้วต่อ BNC & เสียบเข้าหัวโพรบแก้ว & รับสัญญาณความต่างศักย์ระดับมิลลิโวลต์จากหัววัด \\
\bottomrule
\end{tabularx}
\end{table}

\begin{tcolorbox}[enhanced, colback=yellow!10, colframe=red!70!black, arc=3mm, title=\textbf{ข้อควรระวังสำคัญด้านความปลอดภัยของแรงดันไฟฟ้า}]
ชิปประมวลผล ATSAMD51 บนบอร์ด Wio Terminal ทำงานที่ระดับแรงดันลอจิก 3.3 โวลต์เท่านั้น ดังนั้นผู้ปฏิบัติการต้องหมุนปรับทริมพอตตัวล่าง (Offset Trimmer) ของโมดูล pH-4502C ให้จ่ายแรงดันที่สภาวะ pH 7.00 เท่ากับ 1.650 โวลต์เสมอ เพื่อป้องกันไม่ให้แรงดันเกิน 3.3 โวลต์เข้าสู่ขา ADC ของไมโครคอนโทรลเลอร์
\end{tcolorbox}

% ---------------------- บทที่ 4 ----------------------
\chapter{สถาปัตยกรรมปัญญาประดิษฐ์ฝังตัว (TinyML ANN 2-4-1)}

แบบจำลองปัญญาประดิษฐ์ในระบบ iSoil pH xAI เป็นโครงข่ายประสาทเทียมแบบ Multi-Layer Perceptron (MLP) โครงสร้าง 2-4-1 ที่ได้รับการแปลงรูปให้ทำงานบนชิป ARM Cortex-M4F 120MHz ได้อย่างรวดเร็วและใช้หน่วยความจำต่ำมาก

\begin{figure}[htbp]
\centering
\begin{tikzpicture}[scale=0.85, auto, >=stealth]
    \node[draw=rbruNavy, fill=white, circle, inner sep=2mm] (in1) at (0, 2.0) {\small $V$};
    \node[draw=rbruNavy, fill=white, circle, inner sep=2mm] (in2) at (0, -2.0) {\small $T$};

    \node[draw=rbruGold, fill=rbruCodeBg, circle, inner sep=2mm] (h0) at (4, 3.0) {\small $h_0$};
    \node[draw=rbruGold, fill=rbruCodeBg, circle, inner sep=2mm] (h1) at (4, 1.0) {\small $h_1$};
    \node[draw=rbruGold, fill=rbruCodeBg, circle, inner sep=2mm] (h2) at (4, -1.0) {\small $h_2$};
    \node[draw=rbruGold, fill=rbruCodeBg, circle, inner sep=2mm] (h3) at (4, -3.0) {\small $h_3$};

    \node[draw=rbruGreen, fill=rbruGreen!20, circle, inner sep=2mm] (out) at (8, 0) {\small $\text{pH}_{\text{AI}}$};

    \foreach \i in {in1, in2}
        \foreach \j in {h0, h1, h2, h3}
            \draw[->, line width=0.8pt, color=rbruSlate!70] (\i) -- (\j);

    \foreach \j in {h0, h1, h2, h3}
        \draw[->, line width=0.8pt, color=rbruGreen!80] (\j) -- (out);

    \node at (0, 3.2) {\textbf{\color{rbruNavy}Input Layer}};
    \node at (4, 4.0) {\textbf{\color{rbruGold}Hidden Layer (ReLU)}};
    \node at (8, 1.2) {\textbf{\color{rbruGreen}Output Layer}};
\end{tikzpicture}
\caption{สถาปัตยกรรมโครงข่ายประสาทเทียม ANN 2-4-1 ฝังตัวบน Wio Terminal}
\label{fig:ann_arch}
\end{figure}

สมการคณิตศาสตร์ในชั้นซ่อนเร้นใช้ฟังก์ชันกระตุ้นแบบ ReLU $\max(0, z)$ เพื่อขจัดความไม่เป็นเชิงเส้นและลดเวลาการคำนวณในระดับไมโครวินาที:
\begin{equation}
h_i = \max\left(0, b_i + W_{i,0} \cdot V + W_{i,1} \cdot (T - 25.0)\right) \quad (i = 0, 1, 2, 3)
\end{equation}
\begin{equation}
\text{pH}_{\text{AI}} = b_{\text{out}} + \sum_{i=0}^3 (h_i \cdot W_{\text{out}, i})
\end{equation}

% ---------------------- บทที่ 5 ----------------------
\chapter{ระบบฐานข้อมูลและแดชบอร์ดเว็บแอปพลิเคชัน}

เพื่อสนับสนุนการทำเกษตรแม่นยำและการรายงานผลวิจัย ระบบได้เชื่อมโยง Wio Terminal เข้ากับแดชบอร์ดเว็บแอปพลิเคชันและฐานข้อมูลแบบเรียลไทม์:
\begin{enumerate}
    \item \textbf{สะพานส่งข้อมูล (Serial Bridge Daemon)} ใช้สคริปต์ Python (\texttt{tools/serial\_bridge.py}) อ่านข้อมูลจากพอร์ต USB Serial แล้วส่งต่อผ่าน HTTP POST ไปยังเซิร์ฟเวอร์
    \item \textbf{ฐานข้อมูล SQLite3} จัดเก็บประวัติการวัดค่าที่ไฟล์ \texttt{data/soil\_ph.db} ซึ่งมีดัชนีเวลาเพื่อการค้นคืนข้อมูลที่รวดเร็ว
    \item \textbf{เว็บพอร์ทัลความเร็วสูง} พัฒนาด้วย PHP และ JavaScript (Chart.js) แสดงผลกราฟเปรียบเทียบสองมิติ บัตรดัชนีสุขภาพดิน และปุ่มส่งออกข้อมูลเป็น CSV ทันที
\end{enumerate}

\section{ระบบการจัดเก็บไฟล์ข้อมูลการทดลองแบบแยกเซสชันอัตโนมัติ}
เพื่อสนับสนุนการทำวิจัยอย่างมีมาตรฐานและป้องกันการเขียนข้อมูลทับซ้อน ระบบได้รับการออกแบบให้สร้างไฟล์ข้อมูลและเซสชันใหม่ทุกครั้งที่มีการรีเซตหรือเริ่มการทดลองใหม่
\begin{enumerate}
    \item \textbf{การสร้างไฟล์ใหม่บน MicroSD Card อัตโนมัติ (Sequential File Generation)} เมื่อเปิดเครื่องใหม่หรือกดปุ่มรีเซต ตัวประมวลผล Wio Terminal จะตรวจค้นไฟล์ใน MicroSD Card แล้วสร้างไฟล์ลำดับใหม่โดยอัตโนมัติ (\texttt{/exp\_001.csv}, \texttt{/exp\_002.csv}, \texttt{/exp\_003.csv} ...) พร้อมเขียนส่วนหัวคอลัมน์มาตรฐาน 10 คอลัมน์
    \item \textbf{การจัดหมวดหมู่ในฐานข้อมูลตามเซสชัน (Database Session Tracking)} ข้อมูลที่ถูกส่งผ่าน Serial Bridge จะถูกกำกับด้วยรหัส \texttt{session\_id} ตรงกับไฟล์ใน MicroSD Card ทำให้ผู้ใช้งานสามารถเลือกดูประวัติ กราฟ และสถิติย้อนหลังแบบเจาะจงรายรอบการทดลองได้
    \item \textbf{การเริ่มรอบการทดลองใหม่ทันที (On-Demand Session Restart)} ผู้ใช้งานสามารถสั่งเริ่มเซสชันใหม่ได้ 3 ช่องทาง คือ (1) กดปุ่มบนซ้าย (ปุ่ม C) บนตัวเครื่อง Wio Terminal, (2) กดปุ่ม ``เริ่มใหม่ (New)'' บนหน้าเว็บแดชบอร์ด, หรือ (3) ส่งคำสั่ง \texttt{START\_NEW} ผ่านพอร์ต Serial
\end{enumerate}

% ---------------------- บทที่ 6 ----------------------
\chapter{ขั้นตอนปฏิบัติการทดสอบและการใช้งานภาคสนาม}

\section{การเตรียมตัวอย่างดินตามระเบียบวิธีวิจัย}
การเตรียมตัวอย่างดินเพื่อตรวจวัดภาคสนามให้ได้ผลแม่นยำสูงสุด มีขั้นตอนดำเนินการดังนี้
\begin{enumerate}
    \item สุ่มเก็บตัวอย่างดินบนที่ระดับความลึก 0 ถึง 15 เซนติเมตร จากแปลงเกษตร 5 ถึง 10 จุด แล้วนำมารวมกันเป็นตัวแทนของแปลง
    \item ผึ่งดินให้แห้งในที่ร่ม บดให้ละเอียดแล้วร่อนผ่านตะแกรงขนาดตา 2 มิลลิเมตร
    \item ชั่งดิน 20 กรัม ผสมกับน้ำกลั่น 20 มิลลิลิตร (อัตราส่วน 1 ต่อ 1) หรือ 50 มิลลิลิตร (อัตราส่วน 1 ต่อ 2.5) ลงในขวดทดลอง
    \item เขย่าหรือกวนสารละลายเป็นเวลา 30 นาที แล้วตั้งทิ้งไว้ให้ดินตกตะกอน 30 นาที
    \item จุ่มหัวโพรบแก้วลงในสารละลายส่วนใส แล้วกดปุ่มบันทึกค่าลงใน MicroSD Card
\end{enumerate}

\section{เกณฑ์การวินิจฉัยสุขภาพดินเพื่อการเพาะปลูกทุเรียน}
ระบบซอฟต์แวร์บนตัวเครื่องจะทำการประเมินค่าความเป็นกรด-ด่างและแสดงสถานะให้เกษตรกรทราบทันทีบนหน้าจอ LCD ดังนี้
\begin{itemize}
    \item \textbf{pH น้อยกว่า 5.0 (VERY ACIDIC)} ดินมีสภาพกรดรุนแรง ธาตุอาหารหลักถูกตรึง ต้องใส่ปูนโดโลไมต์ปรับสภาพ
    \item \textbf{pH 5.5 ถึง 6.5 (OPTIMAL DURIAN)} ดินมีความสมบูรณ์เหมาะสมที่สุดสำหรับการเจริญเติบโตของทุเรียน
    \item \textbf{pH 6.6 ถึง 7.5 (NEUTRAL SOIL)} ดินมีสภาพเป็นกลาง เหมาะสำหรับพืชไร่และพืชสวนทั่วไป
    \item \textbf{pH มากกว่า 7.5 (ALKALINE SOIL)} ดินเป็นด่าง พืชอาจขาดธาตุสังกะสีและเหล็ก ควรใส่ปุ๋ยอินทรีย์
\end{itemize}

% ---------------------- ประวัติคณะผู้วิจัย ----------------------
\chapter*{ประวัติคณะผู้วิจัย}
\addcontentsline{toc}{chapter}{ประวัติคณะผู้วิจัย}

\begin{tcolorbox}[enhanced, colback=white, colframe=rbruNavy, arc=3mm, boxrule=1pt]
\textbf{\large ผู้ช่วยศาสตราจารย์ ดร.ชีวะ ทัศนา (Co-Principal Investigator)}\\
สาขาวิชาฟิสิกส์ คณะวิทยาศาสตร์และเทคโนโลยี มหาวิทยาลัยราชภัฏรำไพพรรณี\\
\textbf{คุณวุฒิการศึกษา} ปร.ด. (ฟิสิกส์ประยุกต์) สจล. พ.ศ. 2555, วท.ม. (ฟิสิกส์ประยุกต์) สจล. พ.ศ. 2549, วท.บ. (ฟิสิกส์) มน. พ.ศ. 2541\\
\textbf{ความเชี่ยวชาญ} ระบบสมองกลฝังตัว ปัญญาประดิษฐ์และการเรียนรู้เชิงลึก เกษตรฟิสิกส์ดิจิทัล และเซนเซอร์อัจฉริยะ\\
\textbf{ไปรษณีย์อิเล็กทรอนิกส์} chewa.t@rbru.ac.th
\end{tcolorbox}

\vspace{4mm}

\begin{tcolorbox}[enhanced, colback=white, colframe=rbruSlate, arc=3mm, boxrule=1pt]
\textbf{\large อาจารย์ธนพัฒน์ ถิระวุฒิ (Principal Investigator)}\\
หลักสูตร คบ.ฟิสิกส์ คณะวิทยาศาสตร์และเทคโนโลยี มหาวิทยาลัยราชภัฏรำไพพรรณี\\
\textbf{คุณวุฒิการศึกษา} วท.ม. (ฟิสิกส์) จุฬาลงกรณ์มหาวิทยาลัย, วท.บ. (ฟิสิกส์) จุฬาลงกรณ์มหาวิทยาลัย\\
\textbf{ความเชี่ยวชาญ} วัสดุเชิงประกอบ การพัฒนาต้นแบบวิจัย และระบบตรวจวัดทางฟิสิกส์\\
\textbf{ไปรษณีย์อิเล็กทรอนิกส์} tanapat.t@rbru.ac.th
\end{tcolorbox}

\vspace{4mm}

\begin{tcolorbox}[enhanced, colback=white, colframe=rbruGreen, arc=3mm, boxrule=1pt]
\textbf{\large ผู้ช่วยศาสตราจารย์ ดร.นันทพร มูลรังษี (Co-Principal Investigator)}\\
สาขาวิชาเคมี คณะวิทยาศาสตร์และเทคโนโลยี มหาวิทยาลัยราชภัฏรำไพพรรณี\\
\textbf{คุณวุฒิการศึกษา} วท.ด. (เคมี) มหาวิทยาลัยเชียงใหม่, วท.บ. (เคมี) มหาวิทยาลัยเชียงใหม่\\
\textbf{ความเชี่ยวชาญ} เคมีวิเคราะห์ การสร้างเครื่องมือวิเคราะห์ทางเคมีอย่างง่าย และการสอบเทียบมาตรฐาน\\
\textbf{ไปรษณีย์อิเล็กทรอนิกส์} nuntaporn.m@rbru.ac.th
\end{tcolorbox}

\end{document}
